\sectionkicker{Source-backed technical notes}
\subsection*{AI for Science — Agentが実験ワークフローへ入る: Technical Notes}
本欄は記事本文で圧縮した一次資料上の情報を、比較・再検証しやすい形で整理したものである。時系列、確認済みの主張、留意点、一次資料URLを掲載する。
\medskip
\subsection*{Theme at a glance}
\addcontentsline{toc}{subsection}{Theme at a glance}
\begin{center}
\footnotesize
\begin{tabularx}{\linewidth}{@{}p{0.20\linewidth}p{0.10\linewidth}p{0.14\linewidth}X@{}}
\toprule
資料 & 位置づけ & 種別 & 時系列 \\
\midrule
Near-autonomous AI chemist / OAI-M1-03 & 主要資料 & AGENT & 2026-06-17 (研究\_RESULT) \\
GPT-Rosalind June capability update & 補足資料 & モデル更新 & 2026-06-03 (RELEASE) \\
\bottomrule
\end{tabularx}
\end{center}
\normalsize

\begin{technicalnote}{Near-autonomous AI chemist / OAI-M1-03}{主要資料}
\begin{tabularx}{\linewidth}{@{}>{\bfseries}p{0.22\linewidth}X@{}}
組織 & OpenAI \& Molecule.one \\
種別 & AGENT \\
時系列 & 2026-06-17 (研究\_RESULT) \\
\end{tabularx}
\smallskip
{\bfseries 一次資料から整理したtechnical points}
\begin{itemize}[leftmargin=1.5em,itemsep=0.35em]
\item \textbf{Vendor claim}: OpenAI and Molecule.one reported a near-autonomous GPT-5.4-driven chemistry workflow that improved a Chan–Lam reaction; humans selected proposals, operated the lab and validated representative reactions.
\end{itemize}
{\bfseries 読む際の境界}
\begin{itemize}[leftmargin=1.5em,itemsep=0.35em]
\item \textbf{分析上の留意点}: The source is first-party OpenAI material. Capability, benchmark, efficiency, or performance statements remain vendor claims unless explicitly described as externally reproduced.
\end{itemize}
{\bfseries 一次資料}
\begingroup\sloppy
\begin{itemize}[leftmargin=1.5em,itemsep=0.25em]
\item {\scriptsize\url{https://openai.com/index/ai-chemist-improves-reaction}}
\end{itemize}
\endgroup
\end{technicalnote}
\medskip

\begin{technicalnote}{GPT-Rosalind June capability update}{補足資料}
\begin{tabularx}{\linewidth}{@{}>{\bfseries}p{0.22\linewidth}X@{}}
組織 & OpenAI \\
種別 & モデル更新 \\
時系列 & 2026-06-03 (RELEASE) \\
\end{tabularx}
\smallskip
{\bfseries 一次資料から整理したtechnical points}
\begin{itemize}[leftmargin=1.5em,itemsep=0.35em]
\item \textbf{Vendor claim}: OpenAI reported stronger medicinal-chemistry, genomics and lab-work capabilities for GPT-Rosalind, including MedChemBench, GeneBench and LabWorkBench evaluations.
\end{itemize}
{\bfseries 読む際の境界}
\begin{itemize}[leftmargin=1.5em,itemsep=0.35em]
\item \textbf{分析上の留意点}: The source is first-party OpenAI material. Capability, benchmark, efficiency, or performance statements remain vendor claims unless explicitly described as externally reproduced.
\end{itemize}
{\bfseries 一次資料}
\begingroup\sloppy
\begin{itemize}[leftmargin=1.5em,itemsep=0.25em]
\item {\scriptsize\url{https://openai.com/index/introducing-new-capabilities-to-gpt-rosalind}}
\end{itemize}
\endgroup
\end{technicalnote}
\medskip
